% Erratum E2 for inclusion in main.tex.
% Suggested use: % Erratum E2 for inclusion in main.tex.
% Suggested use: % Erratum E2 for inclusion in main.tex.
% Suggested use: % Erratum E2 for inclusion in main.tex.
% Suggested use: \input{erratum_E2_distance_compact_support}

\subsection*{Erratum E2: compact support of the distance function}
\label{s:erratum:E2}

In Section 2.3 (``Distance function''), the sentence
``It is clear that $d$ has compact support if and only if $x \in X_1$''
is not valid for a general Banach space $X_1$ continuously embedded in
$X$.  The correct statement is
\[
	 d \text{ has compact support}
	 \quad\Longleftrightarrow\quad
	 x \in \bigcup_{R>0}
	 \overline{\{x_1\in X_1:\norm{x_1}{1}\leq R\}}^{\,X} .
\]
In particular, $x\in X_1$ implies that $d$ has compact support, but the
converse requires an additional assumption, for instance that the closed
balls of $X_1$ are closed in $X$.

Here is a counterexample to the converse as stated.  Let
$X=\ell^2(\IN_0)$ with orthonormal basis $u_0,u_1,u_2,\ldots$, and let
$Z=\ell^1(\IN)$.  Define a bounded injective operator
$J:Z\to X$ by
\[
	J a
	:=
	\left(\sum_{n\geq 1} a_n\right)u_0
	+
	\sum_{n\geq 1} \frac{a_n}{n} u_n .
\]
Let $X_1:=JZ\subset X$ and equip $X_1$ with the norm
\[
	\norm{J a}{1}:=\norm{a}{\ell^1}.
\]
Then $X_1$ is a Banach space continuously embedded in $X$.  Indeed,
\[
	\norm{Ja}{X}^2
	=
	\left|\sum_{n\geq 1}a_n\right|^2
	+
	\sum_{n\geq 1}\frac{|a_n|^2}{n^2}
	\leq
	2\norm{a}{\ell^1}^2 .
\]
Take $x:=u_0\in X$.  Then $x\notin X_1$: if $Ja=u_0$, then comparison
of the $u_n$-coordinates for $n\geq 1$ gives $a_n=0$ for all $n$, and
then the $u_0$-coordinate is also zero, a contradiction.

On the other hand, for the unit vectors $e_n\in\ell^1(\IN)$ one has
\[
	J e_n = u_0 + \frac1n u_n,
	\qquad
	\Norm{J e_n}{1}=1,
	\qquad
	\Norm{J e_n-u_0}{X}=\frac1n .
\]
Consequently, for the distance function associated with this $x$,
\[
	d(1)
	\leq
	\inf_{n\geq 1}\Norm{J e_n-u_0}{X}
	=0 .
\]
Thus $d(r)=0$ for all $r\geq 1$, so $d$ has compact support, while
$x\notin X_1$.  This disproves the printed equivalence in the stated
generality.

The printed sentence becomes correct under the additional hypothesis that
each ball
\[
	\{x_1\in X_1:\norm{x_1}{1}\leq R\}
\]
is closed in $X$.  Under that hypothesis, $d(R)=0$ implies that $x$ lies
in such a closed ball, hence $x\in X_1$.  This extra hypothesis holds,
for example, when $X_1$ is reflexive and the embedding $X_1\hookrightarrow
X$ is continuous.


\subsection*{Erratum E2: compact support of the distance function}
\label{s:erratum:E2}

In Section 2.3 (``Distance function''), the sentence
``It is clear that $d$ has compact support if and only if $x \in X_1$''
is not valid for a general Banach space $X_1$ continuously embedded in
$X$.  The correct statement is
\[
	 d \text{ has compact support}
	 \quad\Longleftrightarrow\quad
	 x \in \bigcup_{R>0}
	 \overline{\{x_1\in X_1:\norm{x_1}{1}\leq R\}}^{\,X} .
\]
In particular, $x\in X_1$ implies that $d$ has compact support, but the
converse requires an additional assumption, for instance that the closed
balls of $X_1$ are closed in $X$.

Here is a counterexample to the converse as stated.  Let
$X=\ell^2(\IN_0)$ with orthonormal basis $u_0,u_1,u_2,\ldots$, and let
$Z=\ell^1(\IN)$.  Define a bounded injective operator
$J:Z\to X$ by
\[
	J a
	:=
	\left(\sum_{n\geq 1} a_n\right)u_0
	+
	\sum_{n\geq 1} \frac{a_n}{n} u_n .
\]
Let $X_1:=JZ\subset X$ and equip $X_1$ with the norm
\[
	\norm{J a}{1}:=\norm{a}{\ell^1}.
\]
Then $X_1$ is a Banach space continuously embedded in $X$.  Indeed,
\[
	\norm{Ja}{X}^2
	=
	\left|\sum_{n\geq 1}a_n\right|^2
	+
	\sum_{n\geq 1}\frac{|a_n|^2}{n^2}
	\leq
	2\norm{a}{\ell^1}^2 .
\]
Take $x:=u_0\in X$.  Then $x\notin X_1$: if $Ja=u_0$, then comparison
of the $u_n$-coordinates for $n\geq 1$ gives $a_n=0$ for all $n$, and
then the $u_0$-coordinate is also zero, a contradiction.

On the other hand, for the unit vectors $e_n\in\ell^1(\IN)$ one has
\[
	J e_n = u_0 + \frac1n u_n,
	\qquad
	\Norm{J e_n}{1}=1,
	\qquad
	\Norm{J e_n-u_0}{X}=\frac1n .
\]
Consequently, for the distance function associated with this $x$,
\[
	d(1)
	\leq
	\inf_{n\geq 1}\Norm{J e_n-u_0}{X}
	=0 .
\]
Thus $d(r)=0$ for all $r\geq 1$, so $d$ has compact support, while
$x\notin X_1$.  This disproves the printed equivalence in the stated
generality.

The printed sentence becomes correct under the additional hypothesis that
each ball
\[
	\{x_1\in X_1:\norm{x_1}{1}\leq R\}
\]
is closed in $X$.  Under that hypothesis, $d(R)=0$ implies that $x$ lies
in such a closed ball, hence $x\in X_1$.  This extra hypothesis holds,
for example, when $X_1$ is reflexive and the embedding $X_1\hookrightarrow
X$ is continuous.


\subsection*{Erratum E2: compact support of the distance function}
\label{s:erratum:E2}

In Section 2.3 (``Distance function''), the sentence
``It is clear that $d$ has compact support if and only if $x \in X_1$''
is not valid for a general Banach space $X_1$ continuously embedded in
$X$.  The correct statement is
\[
	 d \text{ has compact support}
	 \quad\Longleftrightarrow\quad
	 x \in \bigcup_{R>0}
	 \overline{\{x_1\in X_1:\norm{x_1}{1}\leq R\}}^{\,X} .
\]
In particular, $x\in X_1$ implies that $d$ has compact support, but the
converse requires an additional assumption, for instance that the closed
balls of $X_1$ are closed in $X$.

Here is a counterexample to the converse as stated.  Let
$X=\ell^2(\IN_0)$ with orthonormal basis $u_0,u_1,u_2,\ldots$, and let
$Z=\ell^1(\IN)$.  Define a bounded injective operator
$J:Z\to X$ by
\[
	J a
	:=
	\left(\sum_{n\geq 1} a_n\right)u_0
	+
	\sum_{n\geq 1} \frac{a_n}{n} u_n .
\]
Let $X_1:=JZ\subset X$ and equip $X_1$ with the norm
\[
	\norm{J a}{1}:=\norm{a}{\ell^1}.
\]
Then $X_1$ is a Banach space continuously embedded in $X$.  Indeed,
\[
	\norm{Ja}{X}^2
	=
	\left|\sum_{n\geq 1}a_n\right|^2
	+
	\sum_{n\geq 1}\frac{|a_n|^2}{n^2}
	\leq
	2\norm{a}{\ell^1}^2 .
\]
Take $x:=u_0\in X$.  Then $x\notin X_1$: if $Ja=u_0$, then comparison
of the $u_n$-coordinates for $n\geq 1$ gives $a_n=0$ for all $n$, and
then the $u_0$-coordinate is also zero, a contradiction.

On the other hand, for the unit vectors $e_n\in\ell^1(\IN)$ one has
\[
	J e_n = u_0 + \frac1n u_n,
	\qquad
	\Norm{J e_n}{1}=1,
	\qquad
	\Norm{J e_n-u_0}{X}=\frac1n .
\]
Consequently, for the distance function associated with this $x$,
\[
	d(1)
	\leq
	\inf_{n\geq 1}\Norm{J e_n-u_0}{X}
	=0 .
\]
Thus $d(r)=0$ for all $r\geq 1$, so $d$ has compact support, while
$x\notin X_1$.  This disproves the printed equivalence in the stated
generality.

The printed sentence becomes correct under the additional hypothesis that
each ball
\[
	\{x_1\in X_1:\norm{x_1}{1}\leq R\}
\]
is closed in $X$.  Under that hypothesis, $d(R)=0$ implies that $x$ lies
in such a closed ball, hence $x\in X_1$.  This extra hypothesis holds,
for example, when $X_1$ is reflexive and the embedding $X_1\hookrightarrow
X$ is continuous.


\subsection*{Erratum E2: compact support of the distance function}
\label{s:erratum:E2}

In Section 2.3 (``Distance function''), the sentence
``It is clear that $d$ has compact support if and only if $x \in X_1$''
is not valid for a general Banach space $X_1$ continuously embedded in
$X$.  The correct statement is
\[
	 d \text{ has compact support}
	 \quad\Longleftrightarrow\quad
	 x \in \bigcup_{R>0}
	 \overline{\{x_1\in X_1:\norm{x_1}{1}\leq R\}}^{\,X} .
\]
In particular, $x\in X_1$ implies that $d$ has compact support, but the
converse requires an additional assumption, for instance that the closed
balls of $X_1$ are closed in $X$.

Here is a counterexample to the converse as stated.  Let
$X=\ell^2(\IN_0)$ with orthonormal basis $u_0,u_1,u_2,\ldots$, and let
$Z=\ell^1(\IN)$.  Define a bounded injective operator
$J:Z\to X$ by
\[
	J a
	:=
	\left(\sum_{n\geq 1} a_n\right)u_0
	+
	\sum_{n\geq 1} \frac{a_n}{n} u_n .
\]
Let $X_1:=JZ\subset X$ and equip $X_1$ with the norm
\[
	\norm{J a}{1}:=\norm{a}{\ell^1}.
\]
Then $X_1$ is a Banach space continuously embedded in $X$.  Indeed,
\[
	\norm{Ja}{X}^2
	=
	\left|\sum_{n\geq 1}a_n\right|^2
	+
	\sum_{n\geq 1}\frac{|a_n|^2}{n^2}
	\leq
	2\norm{a}{\ell^1}^2 .
\]
Take $x:=u_0\in X$.  Then $x\notin X_1$: if $Ja=u_0$, then comparison
of the $u_n$-coordinates for $n\geq 1$ gives $a_n=0$ for all $n$, and
then the $u_0$-coordinate is also zero, a contradiction.

On the other hand, for the unit vectors $e_n\in\ell^1(\IN)$ one has
\[
	J e_n = u_0 + \frac1n u_n,
	\qquad
	\Norm{J e_n}{1}=1,
	\qquad
	\Norm{J e_n-u_0}{X}=\frac1n .
\]
Consequently, for the distance function associated with this $x$,
\[
	d(1)
	\leq
	\inf_{n\geq 1}\Norm{J e_n-u_0}{X}
	=0 .
\]
Thus $d(r)=0$ for all $r\geq 1$, so $d$ has compact support, while
$x\notin X_1$.  This disproves the printed equivalence in the stated
generality.

The printed sentence becomes correct under the additional hypothesis that
each ball
\[
	\{x_1\in X_1:\norm{x_1}{1}\leq R\}
\]
is closed in $X$.  Under that hypothesis, $d(R)=0$ implies that $x$ lies
in such a closed ball, hence $x\in X_1$.  This extra hypothesis holds,
for example, when $X_1$ is reflexive and the embedding $X_1\hookrightarrow
X$ is continuous.
